\section{Related Works}
\label{sec:related_work}

%-------------------------------------------------------------------------
\subsection{Positional Encoding for Transformers}

Transformers~\cite{vaswani2017attention} are permutation-invariant and require explicit positional information to model token order or spatial layout. Early methods employ absolute positional embeddings, either fixed sinusoidal or learned, added directly to token features. Vision Transformers~\cite{dosovitskiy2021image} adopt learned 2D absolute embeddings defined on a fixed spatial grid, using bicubic interpolation~\cite{beyer2022better} when transferring to higher resolutions. However, such embeddings are resolution-specific and do not provide true zero-shot extrapolation.

Relative positional embeddings (RPEs) instead encode pairwise displacements by modifying attention scores, ensuring translation invariance and enabling generalization to unseen positions. In NLP, methods such as those of Shaw~\emph{et al.}~\cite{shaw2018self} and T5~\cite{raffel2020exploring} improve sequence-length generalization, while in vision, Swin Transformer~\cite{liu2021swin} introduces learnable 2D relative position biases within local windows. More recent approaches such as FlexiViT~\cite{beyer2023flexivit} and conditional positional encodings~\cite{chu2023conditional} explore resizing strategies for both patch and positional embeddings, but still rely on bounded lookup tables or learned absolute embeddings with limited support for long-range spatial extrapolation.

The challenge of length extrapolation has been most actively studied in large language models. Several key properties have been identified for successful extrapolation~\cite{hong2024token,dong2024exploring}: translation-invariant dependence on relative offsets, monotonic or smoothly decaying attention with distance, and suppression of high-frequency oscillations that destabilize long-range correlations. Building on these principles, a range of methods have been proposed, including RoPE~\cite{su2024roformer}, ALiBi~\cite{press2022train}, YaRN~\cite{peng2024yarn}, FIRE~\cite{li2024functional}, position interpolation~\cite{chen2024extending}, XPOS~\cite{sun2022length}, InfLLM~\cite{xiao2024infllm}, divide-and-conquer scaling~\cite{yang2025dcis}, mixed-radix extensions~\cite{tian2026mrrope}, imaginary RoPE extensions~\cite{liu2025beyond}, multi-scale self-injection~\cite{han2026stacked}, and context window scheduling~\cite{zhu2025skyladder}. While these mechanisms have enabled dramatic context extension in 1D, existing analyses are almost exclusively confined to linear token sequences; how these extrapolation principles transfer to 2D spatial grids, where offsets become vectors and distances grow radially, remains largely unexplored.


\subsection{Positional Embeddings in Vision}

In the visual domain, the adaptation of modern positional encodings is still in its early stages. The original ViT and subsequent works~\cite{dosovitskiy2021image,beyer2022better} interpolate absolute positional embeddings at new resolutions, enabling fine-tuning but not zero-shot extrapolation. Among relative methods, 2D RoPE~\cite{heo2024rope} has been applied to vision transformers by extending rotary embeddings to two spatial axes, but has not been systematically compared against other modern encodings. ALiBi, FIRE, and YaRN, despite their proven 1D extrapolation capabilities, have received little attention in the vision literature.

This gap is our primary motivation. Existing studies evaluate individual methods in isolation or focus exclusively on classification; no prior work provides a controlled comparison of rotation-based and bias-based positional encodings across classification, detection, and segmentation under zero-shot resolution extrapolation. Our study fills this gap by benchmarking five representative methods under identical training and evaluation conditions.
